%! Introducing Vectors — Unit 1, Lesson 1.0 — Homework
\documentclass[10pt]{article}
\usepackage{linalg-article}
\usepackage{linalg-boxes}
\newcommand{\TallMath}[1]{$\displaystyle #1\rule[-1.4em]{0pt}{3.2em}$}
\newcommand{\vv}{\mathbf{v}}
\newcommand{\ww}{\mathbf{w}}

\begin{document}

\pageheader{Unit 1, Lesson 1.0}{Homework}
\namedateperiod

\vspace{0.12in}

\begin{notesbox}{Practice}
Throughout, let $\vv = \begin{bmatrix} 2 \\ 7 \end{bmatrix}$ and
$\ww = \begin{bmatrix} -5 \\ 1 \end{bmatrix}$. Show your steps.

\begin{enumerate}[leftmargin=1.9em, itemsep=12pt]
  \item Compute each vector.
    \begin{enumerate}[label=(\alph*), leftmargin=1.8em, itemsep=4pt]
      \item $\vv + \ww$ \hfill \item $\vv - \ww$ \hfill
      \item $3\vv$ \hfill \item $\vv + 2\ww$
    \end{enumerate}
    \vspace{1.4cm}
  \item Find each length. Leave any non-perfect square in simplest radical form.
    \begin{enumerate}[label=(\alph*), leftmargin=1.8em, itemsep=4pt]
      \item $\left\|\begin{bsmallmatrix} 9 \\ 12 \end{bsmallmatrix}\right\|$ \hfill
      \item $\left\|\begin{bsmallmatrix} -8 \\ 15 \end{bsmallmatrix}\right\|$ \hfill
      \item $\left\|\begin{bsmallmatrix} 1 \\ 1 \end{bsmallmatrix}\right\|$
    \end{enumerate}
    \vspace{1.4cm}
  \item Which of these is a \textbf{unit vector}? Justify with a calculation.
    \[
      \mathbf{p}=\begin{bmatrix} 0 \\ 1 \end{bmatrix}, \qquad
      \mathbf{q}=\begin{bmatrix} 1 \\ 2 \end{bmatrix}
    \]
    \writelines{2}
  \item A hiker walks $\begin{bmatrix} 3 \\ -4 \end{bmatrix}$ then $\begin{bmatrix} 5 \\ 4 \end{bmatrix}$
        (kilometers, east/north).
    \begin{enumerate}[label=(\alph*), leftmargin=1.8em, itemsep=6pt]
      \item What is the hiker's net displacement (the sum)? \vspace{0.8cm}
      \item How far is the hiker from the start, in a straight line? \vspace{0.8cm}
      \item \emph{Interpret:} the hiker walked $5+\sqrt{41}\approx 11.4$ km of trail. Why is
            that more than your answer to (b)? \writelines{2}
    \end{enumerate}
\end{enumerate}
\end{notesbox}

\begin{extensionbox}
A vector can also be a \textbf{data record}. Store a smoothie as
$\begin{bsmallmatrix} \text{cups fruit} \\ \text{cups yogurt} \end{bsmallmatrix}$: one
serving is $\mathbf{s}=\begin{bsmallmatrix} 2 \\ 1 \end{bsmallmatrix}$. Write the vector for a
\emph{triple} batch, and explain what scalar multiplication means in this setting.
\writelines{2}
\end{extensionbox}

\begin{spiralbox}
\textbf{Coming next --- Lesson 1.1: Linear Combinations of Vectors.} Today you scaled
vectors ($c\,\vv$) and added them ($\vv+\ww$). Next we do both at once: $c\,\vv + d\,\ww$.
The big question becomes: \emph{which points in the plane can you reach} by combining two
given vectors?
\end{spiralbox}

\end{document}
